% Modern editorial, markup and figure/layout contributions: CC-BY-SA-4.0.
% Historical text and illustrations: Leonhard Euler (1744), public domain.
% Component scope and upstream attribution: see RIGHTS.md.
\sourcepage{24}
\[
\left.\begin{array}{rcl}
\text{pro }L&=&F_{\prime}\\
\text{pro }K&=&F_{\prime\prime}\\
\text{pro }I&=&F_{\prime\prime\prime}\\
\text{pro }H&=&F_{\mathrm{iv}}\\
&\text{\&c.}
\end{array}\right\}
\quad\text{\textit{pro punctis abscissæ antecedentibus.}}
\]
\textit{Atque hoc pacto, sine prolixa differentialium scriptione, valor functionis cujuscunque variabilis, qui in quovis abscissæ puncto locum obtinet, commode indicabitur.}

\subsection*{Corollarium I.}
\originalsection{49}Cum igitur functionis cujusque valor, in loco quocunque, sit æqualis suo valori in loco antecedente differentiali suo aucto, erit
\[
\begin{array}{rclcrcl}
F'&=&F+dF&&F&=&F_{\prime}+dF_{\prime}\\
F''&=&F'+dF'&&F_{\prime}&=&F_{\prime\prime}+dF_{\prime\prime}\\
F'''&=&F''+dF''&&F_{\prime\prime}&=&F_{\prime\prime\prime}+dF_{\prime\prime\prime}\\
F^{\mathrm{iv}}&=&F'''+dF'''&&F_{\prime\prime\prime}&=&F_{\mathrm{iv}}+dF_{\mathrm{iv}}\\
\multicolumn{3}{c}{\&c.}&&\multicolumn{3}{c}{\&c.}
\end{array}
\]

\subsection*{Corollarium II.}
\originalsection{50}Si ex singulis abscissæ divisionibus applicatæ ducantur, atque ea quæ abscissæ $AM=x$ respondet, nempe $Mm$, ponatur $=y$, reliquæ tam sequentes quam antecedentes, ita denotabuntur:
\[
\begin{array}{rlcrl}
Mm&=y&&Mm&=y\\
Nn&=y'&&Ll&=y_{\prime}\\
Oo&=y''&&Kk&=y_{\prime\prime}\\
Pp&=y'''&&Ii&=y_{\prime\prime\prime}\\
Qq&=y^{\mathrm{iv}}&&Hh&=y_{\mathrm{iv}}\\
\multicolumn{2}{c}{\&c.}&&\multicolumn{2}{c}{\&c.}
\end{array}
\]

\subsection*{Corollarium III.}
\originalsection{51}\sourcepage{25}Cum deinde valor ipsius $p$ sit $=\dfrac{dy}{dx}=\dfrac{Nn-Mm}{dx}$; erit $p=\dfrac{y'-y}{dx}$; sequentes autem pariter ac antecedentes ipsius $p$ valores ita se habebunt:
\[
\renewcommand{\arraystretch}{1.8}
\begin{array}{rclcrcl}
p&=&\dfrac{y'-y}{dx}&&p&=&\dfrac{y'-y}{dx}\\
p'&=&\dfrac{y''-y'}{dx}&&p_{\prime}&=&\dfrac{y-y_{\prime}}{dx}\\
p''&=&\dfrac{y'''-y''}{dx}&&p_{\prime\prime}&=&\dfrac{y_{\prime}-y_{\prime\prime}}{dx}\\
p'''&=&\dfrac{y^{\mathrm{iv}}-y'''}{dx}&&p_{\prime\prime\prime}&=&\dfrac{y_{\prime\prime}-y_{\prime\prime\prime}}{dx}\\
\multicolumn{3}{c}{\&c.}&&\multicolumn{3}{c}{\&c.}
\end{array}
\]

\subsection*{Corollarium IV.}
\originalsection{52}Deinde, quia est $q=\dfrac{dp}{dx}=\dfrac{p'-p}{dx}$ erit
\[
q=\frac{y''-2y'+y}{dx^2},
\]
ex quo quantitatis $q$ valores, cum sequentes tum antecedentes, ita se habebunt:
\[
\renewcommand{\arraystretch}{1.8}
\begin{array}{rclcrcl}
q&=&\dfrac{y''-2y'+y}{dx^2}&&q&=&\dfrac{y''-2y'+y}{dx^2}\\
q'&=&\dfrac{y'''-2y''+y'}{dx^2}&&q_{\prime}&=&\dfrac{y'-2y+y_{\prime}}{dx^2}\\
q''&=&\dfrac{y^{\mathrm{iv}}-2y'''+y''}{dx^2}&&q_{\prime\prime}&=&\dfrac{y-2y_{\prime}+y_{\prime\prime}}{dx^2}\\
\multicolumn{3}{c}{\&c.}&&\multicolumn{3}{c}{\&c.}
\end{array}
\]

\subsection*{Corollarium V.}
\originalsection{53}Simili igitur modo per ista applicatarum signa poterunt \sourcepage{26}valores quantitatum $r$, $s$, $t$, \&c. ut has supra assumsimus, determinari, atque ex figura definiri. Erit scilicet
\[
\begin{aligned}
r&=\frac{y'''-3y''+3y'-y}{dx^3},\\[5pt]
s&=\frac{y^{\mathrm{iv}}-4y'''+6y''-4y'+y}{dx^4},\\[5pt]
t&=\frac{y^{\mathrm{v}}-5y^{\mathrm{iv}}+10y'''-10y''+5y'-y}{dx^5},\\[5pt]
&\text{\&c.}
\end{aligned}
\]
unde harum litterarum valores tam præcedentes quam antecedentes formari possunt.

\Needspace{11\baselineskip}
\subsection*{Corollarium VI.}
\originalsection{54}Quod si autem formula $\int Z\,dx$ ad abscissam $AM=x$ fuerit relata; erit ejus valor sequenti abscissæ elemento $MN=dx$ respondens $=Zdx$. Hincque simili modo formulæ $\int Z\,dx$ valores singulis abscissæ elementis respondentes denotabuntur ut sequitur:
\[
\begin{array}{rclcrcl}
\text{pro }MN&=&Zdx&&\text{pro }MN&=&Zdx\\
\text{pro }NO&=&Z'dx&&\text{pro }LM&=&Z_{\prime}dx\\
\text{pro }OP&=&Z''dx&&\text{pro }KL&=&Z_{\prime\prime}dx\\
\text{pro }PQ&=&Z'''dx&&\text{pro }IK&=&Z_{\prime\prime\prime}dx\\
\multicolumn{3}{c}{\&c.}&&\multicolumn{3}{c}{\&c.}
\end{array}
\]


\subsection*{Corollarium VII.}
\originalsection{55}Si ergo expressio $\int Zdx$ ad abscissam curvæ $AM=x$ pertineat; ejusdem expressionis valor, qui conveniet abscissæ propositæ $AZ$, erit
\[
=\int Zdx+Zdx+Z'dx+Z''dx+Z'''dx+\&c.
\]
in infinitum, donec perveniatur ad ultimum punctum $Z$.

\subsection*{Corollarium VIII.}
\originalsection{56}Si igitur curva inveniri debeat, quæ pro data abscissa $AZ$ \sourcepage{27}valorem formulæ $\int Zdx$ habeat maximum minimumve; tum, posita abscissa quacunque indefinita $AM=x$, efficiendum est ut hæc expressio
\[
\int Zdx+Zdx+Z'dx+Z''dx+Z'''dx+\&c.
\]
usque in $Z$ fiat maxima vel minima.

\subsection*{Scholion.}
\originalsection{57}Quanquam hæc hypothesis tantum pro arbitrio est facta; tamen ista signa maximam afferent utilitatem ad Problemata, quæ ad hanc methodum maximorum \& minimorum pertinent, succincte resolvenda. Plurimum enim valet in hujusmodi negotiis commoda signorum electio, ejusque ope calculus non solum contrahi, sed etiam multo facilior \& expeditior reddi potest. Præstabit autem iste signandi modus longe alteri recepto, quo per differentialia valores functionum variabilium proxime sequentes exprimi solent; eo quod in ipsa resolvendi methodo alius generis differentialia occurrent, quæ cum naturalibus quantitatum variabilium differentialibus facile confundi possent, nisi, ista assumta signandi methodo, naturalia differentialia notatione saltem tollerentur.

\subsection*{Propositio IV. Theorema.}
\originalsection{58}\marginnote{\footnotesize\hyperlink{fig:3}{Fig. 3.}}\textit{Si $amnoz$ fuerit curva ad abscissam datam $AZ$ relata, in qua formula $\int Zdx$ maximum minimumve obtineat valorem; atque alia concipiatur curva $amvoz$ ab ista infinite parum discrepans, tum valor formulæ $\int Zdx$ pro utraque curva erit idem.}

\subsection*{Demonstratio.}
Quando in Analysi formula quæpiam variabilis fit maxima, tum primo crescendo continuo magis ad maximum valorem accedit, deinde vero cum hunc attigit, iterum decrescendo ab eo recedit. Iste autem accessus ad maximum valorem atque recessus ab eodem ita fit, ut dum quantitas proxime ad maximum valorem versatur, tum ejus incrementa ac decrementa \sourcepage{28}momentanea evanescant; hocque idem de minimo est intelligendum. Dantur quidem etiam ejusmodi maxima \& minima, circa quæ incrementa \& decrementa sint infinite magna; verum hujus generis maxima \& minima in præsenti instituto raro locum inveniunt, \& si inveniunt, facile erit ea determinare. Sufficiat igitur notasse circa maximum \& minimum mutationes momentaneas non dari posse finitas. Quod si ergo in curva $amnoz$ expressio $\int Zdx$ maximum minimumve habeat valorem; pro alia curva ejusdem expressionis valor eo magis a maximo minimove recedet, quo magis hæc alia curva ab illa discrepet. Sin autem alia curva infinite parum differat ab illa satisfaciente, tum, pro utraque, formula $\int Zdx$ eundem obtinebit valorem. Hujusmodi autem curvam minime discrepantem concipiemus, si arcum tantum infinite parvum $mno$ infinite parum variari, ejusque loco arcum $mvo$ substitui ponamus. Quamobrem ex curva $az$, pro qua $\int Zdx$ maximum est vel minimum, portionem infinite parvam $mno$ excindi, ejusque loco aliam $mvo$ infinite parum ab illa discrepantem inseri intelligamus; tum valor formulæ $\int Zdx$ qui convenit curvæ $amnoz$ æqualis erit valori, qui convenit curvæ $amvoz$. Q. E. D.

\subsection*{Corollarium I.}
\originalsection{59}Quoniam mutatio debet poni quamminima; non sufficiet arcum $mno$, qui immutari ponitur, accipere infinite parvum, sed etiam deviatio $nv$, præ arcus longitudine $mno$, debet esse infinite parva.

\subsection*{Corollarium II.}
\originalsection{60}Posita igitur tali mutatione in curva, mutatio inde etiam in valore formulæ $\int Zdx$ orietur; quæ autem per demonstrationem erit evanescens. Atque hoc modo ex tali assumta mutatione orietur æquatio, quæ simul curvæ quæsitæ naturam præbebit.

\subsection*{Scholion.}
\originalsection{61}\sourcepage{29}In hac Propositione continetur universa methodus resolvendi Problemata, quibus curva desideratur in qua valor formulæ cujusdam indeterminatæ ut $\int Zdx$ sit maximus vel minimus. Semper enim concipitur portio curvæ infinite parva, uti $mno$, aliquantillum variari in $mvo$, atque tum quæritur differentia valorum quos formula $\int Zdx$, cum pro curva vera $amnoz$, tum pro ficta $amvoz$, sortitur, eaque differentia nihilo æqualis posita dat naturam curvæ quæsitæ. Mutatio autem ista in loco indefinito fieri debet, ut ad totam curvam pertineat, atque ad singula loca pateat. Potest autem ista mutatio utcunque institui, dummodo sit infinite parva, atque vel ad duo vel plura curvæ elementa extendi; semper enim eadem resultare debet æquatio finalis. Interim tamen calculi commoditas postulat, ut mutatio in tam paucis elementis instituatur, quæ sufficiat ad solutionem absolvendam. Ita si, inter omnes omnino curvas eidem abscissæ respondentes, ea determinari debeat, in qua sit $\int Zdx$ maximum vel minimum; tum sufficiet bina tantum curvæ elementa mutata concipere. At si non inter omnes curvas, sed eas tantum quæ unam pluresve expressiones communes habeant, ea definiri debeat in qua quæpiam quantitas sit maxima vel minima; tum mutationem non quamcunque $mvo$ accipere licet, sed talem statui oportet, ut illæ proprietates omnibus curvis communes conserventur. His igitur casibus, duo elementa non sufficient, sed plura accipi debebunt, ut omnibus conditionibus satisfieri queat.

\subsection*{Definitio V.}
\originalsection{62}\textit{Valor Differentialis} datæ maximi minimive formulæ respondens est differentia inter valores, quos hæc formula, cum in ipsa curva quæsita, tum in eadem infinite parum immutata, obtinet.

\subsection*{Corollarium I.}
\originalsection{63}\sourcepage{30}In curva igitur, pro qua data formula, puta $\int Zdx$, maximum minimumve esse debet, hujus formulæ valor differentialis respondens evanescet. Atque hanc ob rem si valor differentialis nihilo æqualis ponatur; habebitur æquatio, qua curvæ quæsitæ natura exprimetur.

\subsection*{Corollarium II.}
\originalsection{64}Ex invento igitur valore differentiali, qui propositæ maximi minimive formulæ respondeat, statim habebitur æquatio exprimens naturam ejus curvæ, in qua formula illa proposita maximum minimumve habeat valorem.

\subsection*{Corollarium III.}
\originalsection{65}Totum igitur negotium ad curvas inveniendas, quæ maximi minimive proprietate gaudeant, eo est reductum, ut pro quaque maximi minimive formula ejus conveniens valor differentialis investigetur.

\subsection*{Scholion.}
\originalsection{66}Cum igitur in genere tradita sit idea non solum naturæ quæstionum, quibus curvæ maximi minimive proprietate præditæ quæruntur, sed etiam methodi, qua ad eas resolvendas uti oporteat, ad ipsam tractationem progrediemur. Ac primo quidem Methodum absolutam, qua curvæ quæruntur quæ inter omnes omnino curvas ad eandem abscissam relatæ maximi minimive proprietate quapiam sint præditæ, trademus. Deinde pergemus ad Methodum maximorum ac minimorum relativam, ad quam tales pertinent quæstiones, quæ non inter omnes curvas datæ abscissæ respondentes, sed eas tantum quæ data quadam communi proprietate una pluribusve gaudent, eam determinari jubent, cui maximi minimive prærogativa quæpiam \sourcepage{31}conveniat. In has autem tractationes natura formulæ $\int Zdx$, quæ maximum minimumve esse debet, ingens discrimen infert, prout $Z$ fuerit functio vel determinata vel indeterminata; quemadmodum jam observavimus.
